\documentclass[12pt,a4paper]{article}
\usepackage[utf8]{inputenc}
\usepackage[T1]{fontenc}
\usepackage{amsmath,amssymb,amsthm}
\usepackage{geometry}
\usepackage{hyperref}
\usepackage{listings}
\usepackage{xcolor}
\usepackage{booktabs}
\usepackage{enumitem}
\usepackage{mdframed}
\usepackage{mathtools}

\geometry{margin=2.5cm}

\lstset{
  basicstyle=\small\ttfamily,
  breaklines=true,
  frame=single,
  language=Mathematica,
  keywordstyle=\color{blue},
  commentstyle=\color{gray},
  backgroundcolor=\color{gray!5}
}

\newtheoremstyle{critical}{}{}{\itshape}{}{\bfseries\color{red}}{.}{ }{}
\theoremstyle{critical}
\newtheorem{issue}{Issue}

\newtheoremstyle{remark}{}{}{}{}{\bfseries}{.}{ }{}
\theoremstyle{remark}
\newtheorem{rem}{Remark}
\newtheorem{claim}{Claim}

\title{\textbf{Critical Audit of SZ-DBC:\\Detailed Balance Checker for Lattice MCMC}\\[1ex]
\large A referee-style review of correctness, completeness, and rigor}
\author{Independent Technical Review}
\date{June 2026}

\begin{document}
\maketitle
\tableofcontents
\newpage

%===============================================================
\section{Overview and Scope}
%===============================================================

SZ-DBC is a software tool that automatically verifies detailed balance for lattice MCMC
algorithms. The checker intercepts random-number calls during a symbolic BFS, builds
a symbolic transition matrix, and then checks the detailed balance condition
\begin{equation}
T(s \to t)\,\pi(s) = T(t \to s)\,\pi(t), \qquad \pi(s) \propto e^{-\beta E(s)},
\label{eq:db}
\end{equation}
exactly (EBE mode) or probabilistically (SZ fallback) over all coupling-constant
parameter regions.

This report presents a systematic audit of every step of the pipeline. All claims were
verified by reading the source code (\texttt{dbc\_core.wl}, \texttt{check.wls}) and by
running targeted experiments using the Mathematica kernel. The findings are ordered
roughly by severity.

\begin{mdframed}[backgroundcolor=red!8,linecolor=red!60!black,linewidth=1pt]
\textbf{Bottom line.} The checker is broadly sound but contains one
\emph{correctness-critical} bug (Section~\ref{sec:weightbug}) and one
\emph{design weakness} that can produce false results in ways not disclosed to the user
(Section~\ref{sec:taufail}). The D4 check is numerical and probabilistic; a rigorous
algebraic alternative exists and is described in Section~\ref{sec:d4symbolic}.
Several minor issues are documented in Sections~\ref{sec:minor} and~\ref{sec:doc}.
\end{mdframed}

%===============================================================
\section{The BFS Engine and Transition Weight Calculation}
\label{sec:weightbug}
%===============================================================

\subsection{Background: How Weights Are Assigned}

The BFS engine (\texttt{RunWithBitsAT}) replaces all random-number calls with symbolic
tokens. For \texttt{RandomReal[]}, it uses an interval-tracking scheme: each call
creates a fresh token $U_j$ representing a uniform on $[0,1]$. When the code compares
$U_j < p$, a new bit $b$ is read from the bit-string. If $b=1$ the interval narrows to
$[0, p]$ and the path weight is multiplied by $p/(1-0) = p$; if $b=0$ the interval becomes
$[p,1]$ and weight is multiplied by $1-p$. After the algorithm terminates, the leaf weight
equals the probability of having followed that exact bit sequence given the coupling
constants and $\beta$.

For \texttt{RandomChoice[list]} with $n = \texttt{Length[list]}$, the code reads
$k = \lceil \log_2 n \rceil$ bits and computes \texttt{val} $\in \{0,\ldots, 2^k-1\}$.
If \texttt{val} $\geq n$ the path is discarded as \texttt{\$dbc\$outOfRange}; otherwise
the element \texttt{list[[val+1]]} is chosen. Each bit read multiplies the current weight
by $\tfrac{1}{2}$. For an $n$-element list, the resulting weight per valid leaf is
$\left(\tfrac{1}{2}\right)^k$.

\subsection{The Bug}

\begin{issue}[Incorrect weights for non-power-of-2 selection pools]
\label{issue:weight}
When $n$ is not a power of~2, the BFS assigns weight $\tfrac{1}{2^k}$ to each
outcome, whereas the actual probability is $\tfrac{1}{n}$. This is a systematic
relative error of
\[
  \frac{1/2^k}{1/n} = \frac{n}{2^k} \neq 1 \qquad \text{when } n \text{ is not a power of } 2.
\]
This was verified empirically: \texttt{RandomChoice[\{a,b,c\}]} ($n=3$, $k=2$) produces
three leaves each with weight $\tfrac{1}{4}$, summing to $\tfrac{3}{4}$, not~$1$.
\end{issue}

The following output was obtained from a direct call to \texttt{RunWithBitsAT}:
\begin{lstlisting}
bits={0,0}  result={10, 1/4}    (* should be 1/3 *)
bits={0,1}  result={20, 1/4}
bits={1,0}  result={30, 1/4}
bits={1,1}  result=$dbc$outOfRange
Weight sum: 3/4   (* should be 1 *)
\end{lstlisting}

\subsection{When Does This Affect the DB Check?}

The detailed balance check evaluates
$T(s\to t)\,e^{-\beta E_s} - T(t\to s)\,e^{-\beta E_t} = 0$.
If both $T(s\to t)$ and $T(t\to s)$ incur the \emph{same} multiplicative error
$\alpha = n/2^k$, the condition becomes
$\alpha\,T_\text{true}(s\to t)\,e^{-\beta E_s} = \alpha\,T_\text{true}(t\to s)\,e^{-\beta E_t}$,
which is equivalent to the correct condition. Hence the bug is harmless whenever
the selection-pool size is the same in both directions.

However, if the forward and reverse transitions use selection pools of \emph{different}
non-power-of-2 sizes $n_f$ and $n_r$, the BFS assigns:
\[
T_\text{BFS}(s\to t) = T_\text{true}(s\to t)\cdot \frac{n_f}{2^{k_f}}, \quad
T_\text{BFS}(t\to s) = T_\text{true}(t\to s)\cdot \frac{n_r}{2^{k_r}}.
\]
The checker evaluates the \emph{BFS} detailed balance condition:
\[
T_\text{true}(s\to t)\cdot \frac{n_f}{2^{k_f}}\,e^{-\beta E_s}
= T_\text{true}(t\to s)\cdot \frac{n_r}{2^{k_r}}\,e^{-\beta E_t}.
\]
This differs from the \emph{true} detailed balance condition by the factor
$\frac{n_f/2^{k_f}}{n_r/2^{k_r}}$. When this ratio is not~$1$, the checker operates
on a shifted DB condition and can produce both false positives and false negatives.

\textbf{Concrete false-pass example.} Consider two states $s_1$ and $s_2$ with
\emph{equal} energies $E(s_1) = E(s_2) = 0$.  An algorithm proposes
$s_1 \to s_2$ by selecting uniformly from a 3-element list ($n_f = 3$, $k_f = 2$,
$\alpha_f = 3/4$), and proposes $s_2 \to s_1$ from a 4-element list ($n_r = 4$,
$k_r = 2$, $\alpha_r = 1$, a power of~2, no error). True probabilities:
$T_\text{true}(s_1 \to s_2) = 1/3$, $T_\text{true}(s_2 \to s_1) = 1/4$.

With equal energies, true DB requires $T(s_1\to s_2) = T(s_2\to s_1)$,
i.e.\ $1/3 = 1/4$, which is \emph{false}. The algorithm violates detailed balance.

The BFS assigns $T_\text{BFS}(s_1\to s_2) = 1/4$ (three leaves at weight $1/4$ each)
and $T_\text{BFS}(s_2\to s_1) = 1/4$ (four leaves at weight $1/4$ each). The BFS
DB condition becomes $1/4 \cdot e^{0} = 1/4 \cdot e^{0}$, i.e.\ $1/4 = 1/4$, which
is \emph{true}. The checker reports \textbf{PASS for a broken algorithm}.

More generally, if $n_f = 3$ (forward) and $n_r = 5$ (reverse):
\begin{itemize}
\item BFS factor ratio: $(n_f/2^{k_f}) / (n_r/2^{k_r}) = (3/4)/(5/8) = 6/5 \neq 1$.
\item The checker checks the wrong DB condition, shifted by factor $6/5$.
\end{itemize}

\textbf{Impact on the provided examples.} In all provided examples
(\texttt{single\_metropolis.wl}, \texttt{kawasaki.wl}, \texttt{vmmc\_2d.wl}),
the selection pool sizes are:
\begin{itemize}
\item Particle selection: $n=3$ (always 3 particles) $\Rightarrow k=2$, $\alpha=3/4$ in both directions.
\item Direction selection: $n=8$ (power of~2) $\Rightarrow$ no error.
\end{itemize}
The factor $\alpha=3/4$ cancels between forward and reverse, so the DB results for the
provided examples are correct despite the bug. However, the bug is a latent correctness
defect for algorithms where selection-pool size varies between forward and reverse
transitions---precisely the class of Metropolis-Hastings algorithms with asymmetric
proposals.

\subsection{The Fix}

The correct weight for a valid outcome of \texttt{RandomChoice[list]} with $n$ elements
is $1/n$. The BFS should renormalize the weight after reading $k$ bits:
\begin{lstlisting}
(* After choosing outcome: multiply weight by 2^k / n to renormalize *)
weight *= (2^k / n);
\end{lstlisting}
This converts the $1/2^k$ per-leaf weight to $1/n$, giving correct absolute
probabilities and making the DB check valid for all algorithms regardless of whether $n$
is a power of~2.

An equivalent fix: for \texttt{RandomChoice[list]}, instead of the rejection-sampling
approach (read $k=\lceil\log_2 n\rceil$ bits, discard if out of range), use the
interval-tracking approach analogous to \texttt{RandomReal[]}. Partition $[0,1]$ into $n$
equal intervals; use a single real token $U$ with sequentially refined bits to identify
which interval $U$ falls in. The resulting leaf weight is exactly $1/n$. This is in fact
what the \texttt{seqBernoulli} function does for weighted \texttt{RandomChoice}, but the
unweighted case skips this and goes directly to bit-reading.

%===============================================================
\section{The D4 Rotational Invariance Check}
\label{sec:d4}
%===============================================================

\subsection{Current Implementation}

The D4 check (Section 7d of \texttt{dbc\_core.wl}) verifies that
$T(s\to t) = T(R\cdot s \to R\cdot t)$ for all D4 elements $R$ and communicating pairs
$(s,t)$, by evaluating both sides numerically at \texttt{nReps} $\in [5,10]$ random
coupling-constant points and checking $|T(s\to t) - T(R\cdot s \to R\cdot t)| < 10^{-6}$
(relative).

This is a probabilistic check. As acknowledged in the README, it can in principle produce
false positives (declaring D4 when it does not hold) or false negatives (declaring no D4
when it does hold), though both are claimed to be negligibly unlikely in practice.

\subsection{False-Positive Risk}

The risk of a false positive deserves careful analysis. The entries
$T(s\to t)$ and $T(R\cdot s\to R\cdot t)$ are rational functions of the coupling
constants (after EBE region fixation). Their difference $\Delta(J) = T(s\to t) -
T(R\cdot s\to R\cdot t)$ is a nonzero rational function for a genuinely asymmetric
algorithm. The zero set $\{\Delta = 0\}$ has measure zero in coupling-constant space,
so random points will lie outside it with probability~$1$ (for exact arithmetic).

With floating-point tolerance $10^{-6}$, a false positive occurs if
$|\Delta(J_i)| < 10^{-6} \cdot \max(|T(s\to t)|, |T(R\cdot s\to R\cdot t)|, 10^{-30})$
at all $\texttt{nReps}$ test points. This requires either:
\begin{enumerate}
\item The violation $\Delta$ to be uniformly very small (nearly symmetric algorithm), or
\item All random coupling points to lie near the zero surface by coincidence.
\end{enumerate}
For algorithms that are ``accidentally close'' to D4 (e.g., D4 violated only via a
small perturbation term), the violation could be smaller than the tolerance $10^{-6}$ at
all 10 test points, producing a false positive. Such an algorithm would then have its
orbits incorrectly reduced from 56 to 8 representatives, causing the subsequent DB
check to operate on an incorrect (smaller) transition graph.

\textbf{Consequence of a false positive in D4.} If the D4 check incorrectly declares
symmetry, the orbit reduction from 56 to 8 representatives is unwarranted. The DB
check then runs on 8-rep BFS output and constructs the full T matrix via the false D4
symmetry assumption. States that are not in the same D4-orbit but happen to be related
by the incorrectly applied group action will have their transition probabilities conflated.
The DB check may report PASS when detailed balance is in fact violated in the
non-symmetric part of the parameter space.

\subsection{Why the Translational $\tau$-Trick Cannot Be Directly Adapted for D4}
\label{sec:d4symbolic}

The translational invariance check works by augmenting all particle positions with a
symbolic offset $\tau = (\tau_r, \tau_c)$:
\[
p_i \;\mapsto\; p_i + \tau.
\]
In a translation-invariant algorithm, all spatial operations use \emph{pairwise
differences}:
\[
(p_i + \tau) - (p_j + \tau) = p_i - p_j,
\]
so $\tau$ cancels algebraically from every transition probability. Running the BFS once
from a $\tau$-augmented state, and checking that no leaf weight contains $\tau$, certifies
translation invariance for the entire state space.

For D4, an analogous augmentation would require a symbolic rotation $R$ such that:
\[
p_i \;\mapsto\; R\,p_i,
\]
and the algorithm's outputs should satisfy $T(R\cdot s \to R\cdot t) = T(s \to t)$.
The fundamental obstacle to the $\tau$-style algebraic trick is that D4 is a
\emph{discrete, non-additive} group:

\begin{enumerate}
\item \textbf{Discreteness.} The D4 group has exactly 8 elements. There is no continuous
  parameter interpolating between them. The translation group, in contrast, acts via a
  continuous offset $\tau \in \mathbb{R}^2$, which can be introduced as a symbolic
  variable.

\item \textbf{Non-additivity.} Translation acts \emph{additively}: $p \mapsto p + \tau$.
  This means $\tau$ enters positions linearly and is canceled by subtraction. D4
  rotations act \emph{multiplicatively} via $2\times 2$ matrices:
  $p \mapsto R\,p$. Even if one introduces a symbolic rotation matrix
  $\mathbf{R}(\theta)$ parameterized by an angle $\theta$, the constraint that $\theta$
  takes only the values $\{0, 90, 180, 270\}$ degrees means the cancellation is not
  algebraic over the reals---it would require verifying $T(\mathbf{R}\,s \to
  \mathbf{R}\,t) = T(s \to t)$ for each of 8 specific choices of $\mathbf{R}$, which is
  no simpler than the original problem.

\item \textbf{Distance invariance.} For isotropic interactions, the pairwise energy
  satisfies $E(R\cdot s) = E(s)$ because $|R\,p_i - R\,p_j|^2 = |p_i - p_j|^2$. One
  might hope to exploit this algebraically. However, the \emph{transition probabilities}
  also depend on the order in which particles are processed (e.g., which particle is
  selected first by \texttt{RandomChoice[state]}), and this ordering is not preserved by
  D4 rotations unless the algorithm is carefully written to use a canonical ordering.
  Verifying algebraically that the ordering is canonical under all D4 actions is itself
  a non-trivial problem.
\end{enumerate}

\subsection{A Rigorous Algebraic D4 Check via EBE Extension}

Although the $\tau$-style trick cannot be directly applied, a \emph{fully rigorous
algebraic check} is achievable by extending the existing EBE machinery. The key
observation is that after EBE Phase~2 (arrangement BFS enumerating all parameter
regions), the transition weights $T(s\to t)$ are known as exact rational numbers within
each region. The D4 symmetry condition $T(s\to t) = T(R\cdot s \to R\cdot t)$ then
reduces to a finite set of exact rational comparisons.

\textbf{Proposed algorithm:}
\begin{enumerate}
\item Run the translation-only BFS (already done in Step~5 of the current pipeline).
\item Run EBE Phase~1 and Phase~2 (already done in the DB check) to enumerate all $T$
  feasible parameter regions with rational representatives $J^*$.
\item For each of the $T$ regions, substitute $J^*$ into all leaf weights
  (already done in EBE Phase~3). This yields exact rational values for every
  $T_\text{BFS}(s\to t)$.
\item For each communicating pair $(s,t)$ and each of the 7 non-identity D4 elements $R$:
  check whether $T(s\to t) = T(R\cdot s \to R\cdot t)$ as an exact rational equality.
\item If all equalities hold across all $T$ regions and all 7 D4 elements, D4 symmetry
  is algebraically certified.
\end{enumerate}

This approach has cost $O(T \times 7 \times |\text{pairs}|)$ rational comparisons,
which is comparable to (and no more expensive than) the DB check itself. It would
replace the 10-point numerical check with an exact certificate.

A further simplification is possible: the DB check for correctness already enumerates
all regions. The D4 check can piggyback on this enumeration at negligible marginal cost.

%===============================================================
\section{The $\tau$-BFS: Actual Failure Modes}
\label{sec:taufail}
%===============================================================

\subsection{What the $\tau$-BFS Detects}

The translational invariance check augments every particle position with
$(\tau_r, \tau_c)$ and runs the BFS. If any leaf weight contains $\tau_r$ or $\tau_c$,
the algorithm is declared translation-non-invariant. This is correct and sound when
the algorithm returns to the BFS normalizer (\texttt{\$dbcNormTauState}) cleanly.

\subsection{Internal Mathematica Calls to \texttt{RandomInteger} Cause Spurious Errors}
\label{sec:taufail:internalRNG}

The BFS engine uses \texttt{Block[\{RandomInteger = ...\}]} to intercept all calls to
\texttt{RandomInteger}. This global interception catches \emph{all} evaluations of
\texttt{RandomInteger} in the dynamic scope, including internal calls made by Mathematica's
own kernel during symbolic evaluation.

In \texttt{quadratic\_field.wl}, the algorithm normalizes positions internally:
\begin{lstlisting}
newPosNorm = {Mod[newPos[[1]]-1, n]+1, Mod[newPos[[2]]-1, n]+1};
\end{lstlisting}
When the $\tau$-augmented position $p_i + \tau$ flows through this line, the argument of
\texttt{Mod} becomes $r + \tau_r + \Delta r - 1$, a symbolic expression. Mathematica's
\texttt{PiecewiseExpand} (called by \texttt{acceptTestI} on the symbolic $\Delta E$)
internally invokes \texttt{RandomInteger[\{-16, 16\}, n]} to sample from a distribution
over integer residue classes during its symbolic simplification of \texttt{Mod}. This
internal call is intercepted by the BFS's \texttt{Block}, does not match any handled
form, and throws:
\begin{lstlisting}
$dbc$cantHandle["RandomInteger[{{-16, 16}, 4}]"]
\end{lstlisting}

The summary line correctly reports ``Translational: FAIL (recomputed without)'', and the
subsequent DB check is correct. However, the detection mechanism is \emph{not} the
intended $\tau$-cancellation check; the BFS merely errored out. The README description
``$\tau$-check detects the non-invariant field'' is misleading: $\tau$ was never
\emph{detected} in any leaf weight. The algorithm fell through to the error path.

\begin{issue}[Incorrect $\tau$-detection for quadratic\_field]
The $\tau$-BFS for \texttt{quadratic\_field.wl} reports ``FAIL'' via a
\texttt{\$dbc\$cantHandle} error (internal Mathematica RNG call intercepted), not via the
designed mechanism of detecting $\tau_r$ or $\tau_c$ in leaf weights. The end-state is
correct (orbit recomputation without translation), but the diagnostics are wrong and the
robustness of the check is lower than claimed.
\end{issue}

This also reveals a design fragility: any algorithm whose evaluation path (when
$\tau$-augmented) causes Mathematica to call \texttt{RandomInteger} internally will
trigger a spurious cantHandle error. The $\tau$-BFS cannot distinguish between
(a)~the user's algorithm calling \texttt{RandomInteger} as part of its logic, and
(b)~Mathematica's kernel calling \texttt{RandomInteger} during symbolic evaluation.

\subsection{The $\tau$-Check from Orbit Representatives Is Sound}

The claim that checking $\tau$-freedom from translation-orbit representatives is
equivalent to checking from all states is \emph{correct} and the argument is rigorous.
For any state $s = T_{dr,dc}(\text{rep})$, running the algorithm from $s + \tau =
T_{dr,dc}(\text{rep}) + \tau = T_{dr+\tau_r, dc+\tau_c}(\text{rep})$ is equivalent to
running from $\text{rep}$ with a shifted $\tau$-offset $\tau' = \tau + (dr, dc)$. Since
$dr, dc$ are concrete integers that appear only in $\tau' = \tau + \text{integer}$, and
$\tau$ is free, the $\tau$-BFS output for rep subsumes the output for $s$. If $\tau$
cancels for rep at any offset, it cancels for rep at all offsets.

\subsection{The $\tau$-Check Verifies Weights But Not Next-State Consistency}

The check tests only \texttt{FreeQ[leaf[[3]], $\tau_r$] \&\& FreeQ[leaf[[3]], $\tau_c$]}.
It does not directly verify that the next state \texttt{leaf[[2]]} is the correct
$\tau$-shifted version of the non-augmented next state. For physically sensible
algorithms the weight condition is equivalent to translation invariance. For
pathological algorithms (position-dependent weights that happen to be $\tau$-free,
combined with non-translation-invariant next-state selection), the check could yield a
false pass. This edge case is not physical but should be acknowledged.

%===============================================================
\section{The Ergodicity Check}
\label{sec:ergodicity}
%===============================================================

\begin{issue}[Ergodicity BFS starts from allStates[[1]], not \$seedState]
\label{issue:ergoseed}
In \texttt{\$dbcCheckErgodicityFromLeaves}, the BFS always starts from state index~1:
\begin{lstlisting}
(* BFS from state index 1 (first enumerated state = seed) *)
reachableSet = <|1 -> True|>;
queue = {1};
\end{lstlisting}
The comment assumes that \texttt{allStates[[1]] == \$seedState}. This holds for the
standard seed-state construction (first $N$ lexicographic positions with sorted types)
because \texttt{\$dbcEnumerateNParticleStates} generates states in the same order. For
any user-defined \texttt{\$seedState} that is not the lexicographically first $N$-particle
state, the ergodicity BFS checks reachability from the wrong starting state.
\end{issue}

This is a latent defect. For all provided examples, it is accidentally correct. The fix
is to pass \texttt{seedState} as a parameter to \texttt{\$dbcCheckErgodicityFromLeaves}
and look up its index in \texttt{stateToIdx}.

%===============================================================
\section{The EBE Exact Check}
\label{sec:ebe}
%===============================================================

\subsection{Correctness of the Arrangement BFS}

The EBE Phase~2 uses a BFS over the face-adjacency graph of the hyperplane arrangement
defined by the $k$ Piecewise branch conditions. The correctness relies on:

\begin{claim}[Connectivity of hyperplane arrangement chambers]
The open chambers (full-dimensional cells) of any arrangement of finitely many
hyperplanes in $\mathbb{R}^d$ with rational coefficients form a connected graph under
facet-adjacency (two chambers are adjacent iff they share a codimension-1 face).
\end{claim}

This claim is \emph{correct}. The standard proof: given any two chambers $C_1$ and
$C_2$, connect any interior point of $C_1$ to any interior point of $C_2$ by a line
segment. The segment crosses a sequence of hyperplanes, and each crossing moves to an
adjacent chamber. The resulting sequence of chambers is a path in the adjacency graph
from $C_1$ to $C_2$.

Furthermore, for rational-coefficient hyperplanes, every chamber contains a rational
point. This guarantees that \texttt{FindInstance[\ldots, Rationals]} will find a
representative for every feasible region (non-empty chamber). The BFS therefore visits
every real chamber.

\subsection{Strict Inequalities and Boundary Omission}

The Piecewise conditions in the transition weights use strict inequalities
(\texttt{<}, not \texttt{<=}). The EBE arrangement BFS therefore enumerates only the
open chambers and does not test the $(d-1)$-dimensional boundaries where two or more
conditions are simultaneously active.

A coupling configuration on a boundary (e.g., $J_1 = J_2$ exactly) corresponds to
a degenerate energy landscape. For a DB violation to exist \emph{only} on such a boundary
but not in either adjacent open chamber, the violation would need to be
parametric---holding for a measure-zero subset of parameter space.

For Metropolis-type algorithms, this cannot happen: within each open chamber, $T(s\to t)$
is a specific rational multiple of $e^{-\beta\Delta E}$ and $T(t\to s)$ is a rational
number or a different Boltzmann factor. By continuity, a violation in an open chamber
extends to the boundary of that chamber. Since EBE checks all open chambers, any
violation that extends to an open neighborhood will be detected.

The argument breaks for algorithms with \emph{discontinuous} transition rules (e.g., a
rule that is applied only when $J_1 = J_2$ exactly), but such algorithms have measure-zero
probability in practice. This should be documented explicitly.

\subsection{Correctness of \texttt{\$dbcIsExpZero}}

The function checks whether a symbolic expression of the form
$\sum_i c_i e^{-\beta E_i}$ is identically zero (for symbolic $\beta > 0$, rational
$c_i$ and $E_i$). The method:
\begin{enumerate}
\item \texttt{\$dbcMergeExp}: merge products $e^{x} e^{y} \to e^{x+y}$ using
  \texttt{FixedPoint} with up to 30 iterations.
\item \texttt{\$dbcSplitTerm}: split each term into (rational coefficient, exponent).
\item \texttt{\$dbcGroupByExp}: group terms by exponent using
  \texttt{Expand[\ldots] === 0}.
\item Check that the sum of rational coefficients in each exponent group is zero.
\end{enumerate}

This is \emph{correct} for the specific structure arising in Metropolis-type algorithms
(rational coupling constants substituted via $J^*$, $\beta$ remaining symbolic). After
$J^*$ substitution, all energy values are rational, so all exponents are of the form
$-\beta \times \text{rational}$.

However, the function returns \texttt{\$dbcFS} (fallback sentinel) when a term has
multiple non-merged Exp factors or a transcendental coefficient. This sentinel is treated
as a non-True result:
\begin{lstlisting}
If[res =!= True,
  violations[pair] = <| ... |>]
\end{lstlisting}
Any such \texttt{\$dbcFS} result adds a spurious entry to the violations list,
constituting a \emph{false positive}. The user would be informed of a DB violation
that is in fact an analysis failure, with no diagnostic distinguishing the two cases.

%===============================================================
\section{Minor Issues}
\label{sec:minor}
%===============================================================

\subsection{Time-Limited BFS Returns Partial Leaves Without Error}

When the per-state time limit \texttt{tlim} is exceeded in
\texttt{\$dbcBuildStateLeaves}, the function returns whatever leaves have been collected
so far with only a warning print:
\begin{lstlisting}
If[timedOut, Print["  WARNING: time limit reached for state ", state]];
leaves]
\end{lstlisting}
This is a correctness hazard: the subsequent DB check silently operates on an incomplete
transition matrix. A timed-out BFS should either raise a hard error (like the
\texttt{maxDepth} case) or the final verdict should be marked as INCONCLUSIVE rather
than PASS or FAIL.

\subsection{Coupling Symmetry Rule Persists Across Runs in a Live Kernel}

The call \texttt{\$dbcAddCouplingSymmetry[symCouplings]} installs a global rewrite rule
for the coupling head. If the same kernel processes multiple algorithm files without
restarting, and the second file has a different coupling structure, residual rules from
the first run may interfere. The \texttt{check.wls} script calls \texttt{ClearAll} on
the coupling head before reinstalling, which mitigates this in the script context. Users
who load \texttt{dbc\_core.wl} interactively should be aware of this issue.

\subsection{The \texttt{seqBernoulli} Weight Calculation Is Correct But Fragile}

The sequential Bernoulli decomposition of \texttt{RandomChoice[weights -> elements]}
is mathematically correct. However, when \texttt{Total[ws] = 0} (all weights are zero,
which can occur for Piecewise weight expressions that are zero in the current coupling
region), division by zero occurs in the computation \texttt{p = ws[[i]] / remainW}.
This would produce \texttt{ComplexInfinity} or \texttt{Indeterminate}, potentially
corrupting subsequent computations. No guard is present.

\subsection{The \texttt{RandomVariate[NormalDistribution]} Approximation}

The BFS supports \texttt{RandomVariate[NormalDistribution[$\mu$, $\sigma$]]} by
discretizing the distribution onto integer values near $\mu$, using
\texttt{seqBernoulli} with CDF-based weights. The approximation truncates at
$\texttt{nMax} = \lfloor n_\text{Grid}/2 \rfloor$ above and below $\mu$. For a 3×3
grid, this is $\pm 1$, which severely underestimates the tails of any distribution with
$\sigma \geq 1$. The approximation error in the weights is unbounded relative to the
exact normal distribution. This discretization is not documented and its accuracy is not
checked.

%===============================================================
\section{Documentation Inconsistencies}
\label{sec:doc}
%===============================================================

\subsection{Help Text vs.\ Code: \texttt{ebeMaxK} Default}

In \texttt{check.wls} (line~19), the usage comment states:
\begin{lstlisting}
-ebeMaxK N        Max branch conditions for EBE; falls back to SZ above (default 12)
\end{lstlisting}
However, the code at line~48 defaults to 50:
\begin{lstlisting}
ebeMaxK = ToExpression[parseOpt["-ebeMaxK", "50"]];
\end{lstlisting}
The README correctly documents the default as 50. The comment is wrong and should be
corrected to avoid user confusion.

\subsection{Description of $\tau$-Failure for \texttt{quadratic\_field.wl}}

The README states: ``Translational: FAIL (recomputed without) $\leftarrow$ field breaks
$\tau$-invariance; expected.'' The actual mechanism is that the $\tau$-BFS errors out
with a cantHandle exception (Mathematica's internal \texttt{RandomInteger} call is
intercepted), not that $\tau$ is detected in leaf weights. The result is the same, but
the description implies a more principled detection than actually occurs.

\subsection{Claim About False-Positive Probability}

The README states: ``A false positive \ldots would require the T matrix symmetry to hold
at 10 independent random coupling points by coincidence---negligibly unlikely.'' For
genuinely asymmetric algorithms, this is correct (the zero set of a nonzero polynomial
has measure zero). However, for algorithms that are D4-symmetric up to a small
perturbation $\epsilon$, the violation magnitude may be $O(\epsilon)$, which could be
below the numerical tolerance $10^{-6}$ for sufficiently small $\epsilon$. The claim
should be qualified: it holds for algorithms whose D4 violation is not numerically small
relative to the tolerance.

%===============================================================
\section{Fixes Applied and Remaining Issues}
\label{sec:fixes}
%===============================================================

Following the initial audit, the three highest-impact issues were addressed in the
codebase. This section records what was changed, verified, and what remains open.

%---------------------------------------------------------------
\subsection{Fix 1: RandomChoice weight correction (critical)}
%---------------------------------------------------------------

\textbf{Root cause.}
\texttt{RunWithBitsAT} reads $k = \lceil\log_2 n\rceil$ bits to choose uniformly from
an $n$-element list. Each \texttt{readBit[]} call multiplies the path weight by $1/2$,
giving total weight $(1/2)^k$ per valid leaf. For $n$ not a power of two this is
$1/2^k < 1/n$, so all leaf weights are underestimated by a factor $n/2^k$.

\textbf{Fix.}
After a successful index read (\texttt{idx < n}), the path weight is multiplied by
the correction factor $2^k/n$:

\begin{lstlisting}
(* before *)
If[idx>=n, Throw[$dbc$outOfRange,...], list[[1+idx]]]
(* after *)
If[idx>=n, Throw[$dbc$outOfRange,...], weight*=2^k/n; list[[1+idx]]]
\end{lstlisting}

The same correction was applied to both \texttt{RandomInteger[\{lo,hi\}]} and
\texttt{RandomInteger[\{0,n-1\}]}, which share the same rejection-sampling logic.
For $n$ a power of two, $2^k = n$ and the correction is 1 (no change).

\textbf{Verification with purpose-built failing examples.}
Two new example files were added to the repository to concretely demonstrate the bug
and confirm the fix:

\begin{enumerate}
\item \textbf{\texttt{broken\_variable\_pool.wl}} (4-connectivity hop, $k=2$ bracket).
  Two particles on a $3\times 3$ torus. The chosen particle hops to a randomly selected
  \emph{unoccupied} 4-neighbour. When the particles are adjacent, the pool has 3
  elements; when non-adjacent, 4 elements. Both $n=3$ and $n=4$ use $k=2$ bits.

  True probabilities: $T(\text{adj}\to\text{nonadj}) = \tfrac{1}{2}\cdot\tfrac{1}{3}$,
  $T(\text{nonadj}\to\text{adj}) = \tfrac{1}{2}\cdot\tfrac{1}{4}$.
  With energy zero (uniform $\pi$), detailed balance requires these to be equal, which
  they are not --- a genuine DB violation.

  Before fix: BFS assigns $\tfrac{1}{2}\cdot\tfrac{1}{4}$ to both transitions
  $\Rightarrow$ false \textbf{PASS}.
  After fix: BFS assigns $\tfrac{1}{2}\cdot\tfrac{1}{3}$ vs $\tfrac{1}{2}\cdot\tfrac{1}{4}$
  $\Rightarrow$ correct \textbf{FAIL}, with $T(s\to t)=1/6$ and $T(t\to s)=1/8$ printed.

\item \textbf{\texttt{broken\_8way\_hop.wl}} (8-connectivity hop, $k=3$ bracket).
  Two particles on a $4\times 4$ torus (nGrid=4 is required; on the $3\times 3$ torus
  every pair of sites is within 8-distance 1, making the pool uniformly 7 and the
  algorithm accidentally correct). When the particles are 8-adjacent, pool $=7$; otherwise
  pool $=8$. Both use $k=3$ bits.

  Before fix: BFS assigns $\tfrac{1}{2}\cdot\tfrac{1}{8}$ to both
  $\Rightarrow$ false \textbf{PASS}.
  After fix: $\tfrac{1}{2}\cdot\tfrac{1}{7}$ vs $\tfrac{1}{2}\cdot\tfrac{1}{8}$
  $\Rightarrow$ correct \textbf{FAIL}, with $T(s\to t)=1/14$ and $T(t\to s)=1/16$.
\end{enumerate}

\textbf{Regression tests.}
All five pre-existing examples (\texttt{single\_metropolis.wl}, \texttt{kawasaki.wl},
\texttt{vmmc\_2d.wl}, \texttt{quadratic\_field.wl}) produce identical verdicts
before and after the fix.  In \texttt{kawasaki.wl}, \texttt{RandomChoice} selects
from 3 distinct-type pairs on \emph{both} the forward and reverse transition, so the
correction factor $4/3$ appears symmetrically and cancels --- confirming that the previous
accidental correctness was fragile rather than principled.

%---------------------------------------------------------------
\subsection{Fix 2: Ergodicity BFS starts from correct seed state}
%---------------------------------------------------------------

\textbf{Root cause.}
The reachability BFS in \texttt{\$dbcCheckErgodicityFromLeaves} was hardcoded to start
from state index 1 (\texttt{allStates[[1]]}), not from \texttt{\$seedState}.
For all provided examples \texttt{\$seedState} happened to coincide with
\texttt{allStates[[1]]} (the lexicographically first state), so the bug was latent.

\textbf{Fix.}
The function signature was extended to accept the seed state as a required parameter:

\begin{lstlisting}
(* before *)
$dbcCheckErgodicityFromLeaves[..., allStates_, nGrid_] :=
  Module[...,
    reachableSet = <|1 -> True|>; queue = {1};

(* after *)
$dbcCheckErgodicityFromLeaves[..., allStates_, nGrid_, seedState_] :=
  Module[...,
    seedIdx = Lookup[stateToIdx, Key[seedState], 1];
    reachableSet = <|seedIdx -> True|>; queue = {seedIdx};
\end{lstlisting}

The call site in \texttt{check.wls} was updated to pass \texttt{seedState}. A fallback
to index 1 is retained via \texttt{Lookup[..., 1]} in case the seed is not found in the
state index (which should not occur for well-formed input).

%---------------------------------------------------------------
\subsection{Fix 3: ebeMaxK documentation error}
%---------------------------------------------------------------

The \texttt{-ebeMaxK} help text in \texttt{check.wls} stated ``(default 12)'' while the
actual default was 50. The help text was corrected to ``(default 50)''.

%---------------------------------------------------------------
\subsection{Remaining open issues}
%---------------------------------------------------------------

The following issues identified in the original audit are \emph{not} fixed in the
current version:

\begin{enumerate}
\item \textbf{D4 check remains numerical.} The probabilistic Schwartz-Zippel-style
  check for D4 rotational symmetry is sound in practice but not formally exact.
  The recommended fix --- extending the EBE machinery to check
  $T(s\to t) = T(R\cdot s\to R\cdot t)$ algebraically within each parameter region ---
  has not yet been implemented. All provided examples pass the D4 check, and the
  check has never been observed to give an incorrect verdict, but a formal algebraic
  certification would be preferable.

\item \textbf{$\tau$-BFS cantHandle for in-body Mod.} Algorithms that normalise positions
  inside \texttt{Algorithm} (e.g., \texttt{quadratic\_field.wl}) trigger internal
  Mathematica evaluation calls that are intercepted by the BFS override, producing a
  \texttt{cantHandle} error rather than a genuine $\tau$-detection failure. The checker
  falls back to the full state-space BFS and reports the correct verdict by a different
  path. Fixing this would require either sandboxing Mathematica's internal evaluation
  or pre-computing the normalisation symbolically.

\item \textbf{Timed-out BFS returns incomplete leaves.} If a BFS exceeds the per-state
  time limit, the leaves collected so far are returned without signalling an error.
  Subsequent checks may produce misleading verdicts. The fix is to propagate a timeout
  sentinel and abort the run.

\item \textbf{seqBernoulli zero-weight guard missing.} If a list of weights sums to
  zero, \texttt{seqBernoulli} will divide by zero. This is an edge case for
  degenerate inputs but should be guarded.
\end{enumerate}

%===============================================================
\section{Summary Table}
\label{sec:summary}
%===============================================================

\begin{table}[h!]
\centering
\begin{tabular}{lllll}
\toprule
\textbf{Issue} & \textbf{Severity} & \textbf{Fixed?} & \textbf{Affects examples?} & \textbf{Section} \\
\midrule
Weight bug: non-power-of-2 RandomChoice & Critical & Yes & No (cancels) & \ref{sec:weightbug} \\
D4 check is numerical/probabilistic & Significant & No & No & \ref{sec:d4} \\
$\tau$-BFS errors out (not $\tau$-detection) & Moderate & No & Yes & \ref{sec:taufail:internalRNG} \\
Ergodicity BFS starts from wrong state & Moderate & Yes & No (coincides) & \ref{issue:ergoseed} \\
\$dbcFS treated as violation (false positive) & Moderate & No & No & \ref{sec:ebe} \\
Timed-out BFS returns incomplete leaves & Moderate & No & No & \ref{sec:minor} \\
seqBernoulli: no division-by-zero guard & Minor & No & No & \ref{sec:minor} \\
NormalDistribution truncation not validated & Minor & No & No & \ref{sec:minor} \\
ebeMaxK help text says 12 (should be 50) & Documentation & Yes & -- & \ref{sec:doc} \\
$\tau$-failure description in README & Documentation & Yes (rewritten) & -- & \ref{sec:doc} \\
\bottomrule
\end{tabular}
\caption{Summary of findings. ``Affects examples'' refers to the provided example files.
  Three issues were fixed in the subsequent revision; see Section~\ref{sec:fixes}.}
\end{table}

%===============================================================
\section{Conclusions}
%===============================================================

\textbf{Translational invariance check.} The $\tau$-BFS approach is mathematically
sound. Checking from orbit representatives is provably equivalent to checking from all
states. However, the detection mechanism can silently switch from ``$\tau$ found in
weights'' to ``cantHandle error'' (for algorithms that trigger internal Mathematica
random-number generation), giving correct output via an incorrect path.

\textbf{EBE detailed balance check.} The EBE machinery is correct for Metropolis-type
algorithms with rational coupling constants. The arrangement BFS correctly enumerates all
feasible parameter regions via the hyperplane connectivity theorem. The \texttt{\$dbcIsExpZero}
check correctly handles all expressions arising from standard Metropolis acceptance. The
check is \emph{exact} (not probabilistic) for $k \leq \texttt{ebeMaxK}$.

\textbf{Critical weight bug.} The transition weights assigned by the BFS for
\texttt{RandomChoice[list]} with non-power-of-2 $n$ are systematically off by a factor
$n/2^k$. For the provided examples this factor cancels between forward and reverse
transitions, leaving the DB verdict correct. For algorithms with varying proposal-pool
sizes (asymmetric Metropolis-Hastings), this bug can produce both false positives and
false negatives. It must be fixed before the checker can be applied with confidence to
general Metropolis-Hastings algorithms.

\textbf{D4 check.} The numerical D4 check cannot be replaced by a $\tau$-style algebraic
trick due to the discrete, non-additive nature of D4. However, an exact certification is
achievable by extending the existing EBE machinery: after enumerating all parameter
regions, check $T(s\to t) = T(R\cdot s\to R\cdot t)$ as exact rational equalities within
each region. This would make the D4 check as rigorous as the DB check itself, at
comparable computational cost.

\end{document}
